% Book 4, Year 2, chapter 22 -- Regular expressions and the limits of regular languages.
% Provenance (never printed): sources/book4/BRIEFS.md, chapter 22.
\chapter{Regular Expressions and the Limits of Regular Languages}\label{ch:b2:regular-languages}

A search box accepts a pattern: \code{(a|a)*b} asks for some $a$'s, each written in one of two ways, then a $b$. On a
text of $a$'s with no $b$, Python's pattern engine answers ``no match''. The program times that answer:

\omcode{code/bachelor-2/22-regular-languages/listings/backtracking.py}{1}{25}{Timing one pattern on longer and longer
texts, each run in its own bounded \omterm{def:b2:processes-and-system-calls:process}{process}}\label{lst:b2:regular-languages:hook}

\omoutput{code/bachelor-2/22-regular-languages/outputs/backtracking.txt}{Python's \code{re} on \code{(a|a)*b} and
$a^n$, on the laptop.}\label{out:b2:regular-languages:hook}

Each extra $a$ doubles the time: 27 letters take more than six seconds, and a search box that receives such a pattern
freezes. The pattern is not hard. It means the same as \code{a*b}, and the previous chapter's automata decide it in one
step per letter. This chapter connects the two: patterns, called \omterm{def:b2:regular-languages:regex}{regular expressions}, describe exactly the languages
that finite automata recognise; every automaton has a smallest equivalent one, which is unique; and some simple
languages, such as balanced brackets, are beyond both. It ends with a pattern engine that cannot freeze.

\begin{recall}
\Cref{ch:b2:finite-automata}: words, languages, concatenation $LM$ and star $L^*$; deterministic automata (DFA),
complete and trim; nondeterministic automata (NFA) with \omterm{def:b2:finite-automata:epsilon}{spontaneous transitions}, which the \omterm{met:b2:finite-automata:subset}{subset construction} turns
into DFAs, at a cost of up to $2^k$ states; the \omterm{def:b2:finite-automata:product}{product automaton}, so that \omterm{def:b2:finite-automata:recognisable}{recognisable languages} are closed under
union, intersection and complement. The Year 1 volume: structural induction on trees, proving a property for leaves
and then for each node from its children.
\end{recall}

\section{Regular expressions}\label{sec:b2:regular-languages:regex}

\begin{definition}[Regular expression, denoted language]\label{def:b2:regular-languages:regex}
The \emph{regular expressions}\index{regular expression} over an \omterm{def:b2:finite-automata:word}{alphabet} $\Sigma$ are the trees built by these rules:
$\emptyset$ and $\varepsilon$ are expressions, every letter $a \in \Sigma$ is an expression, and if $e$ and $f$ are
expressions, so are $e \mid f$ (alternation), $ef$ (concatenation) and $e^*$ (star). The \emph{language
denoted}\index{denoted language} by an expression is defined on the tree: $L(\emptyset) = \emptyset$, $L(\varepsilon)
= \{\varepsilon\}$, $L(a) = \{a\}$, $L(e \mid f) = L(e) \cup L(f)$, $L(ef) = L(e) L(f)$ and $L(e^*) = L(e)^*$. Two
expressions are \emph{equivalent} if they denote the same language.
\end{definition}

Written on one line, star binds tighter than concatenation, which binds tighter than alternation: $ab^* \mid c$ is
$(a(b^*)) \mid c$. The \emph{size} of an expression is the number of nodes of its tree. So $(a \mid b)^*abb$ denotes the
words over $\{a, b\}$ that end with $abb$, and $(ab \mid b)^*$ the words in which every $a$ is followed by a $b$. The
library of this chapter's listings, \code{regex.rs}, holds the tree as a Rust enumeration with six cases, and its
function \code{lang} is the definition run, cut at a length $n$ so that the sets stay finite:

\omcode{code/bachelor-2/22-regular-languages/listings/regex.rs}{159}{186}{The denotation, by induction on the tree;
the star adds one more factor until nothing new fits}\label{lst:b2:regular-languages:lang}

It is slow, since it builds every word, but it is obviously right, and the tests check the fast code against it.

\begin{definition}[Regular language]\label{def:b2:regular-languages:regular}
A language is \emph{regular}\index{regular language} if some \omterm{def:b2:regular-languages:regex}{regular expression} denotes it. Equivalently, the regular
languages form the smallest class of languages that contains $\emptyset$, $\{\varepsilon\}$ and every $\{a\}$, and is
closed under union, concatenation and star.
\end{definition}

The two descriptions agree by structural induction: every \omterm{def:b2:regular-languages:regex}{denoted language} is in any such class, and the \omterm{def:b2:regular-languages:regex}{denoted
languages} form such a class. Every finite language is \omterm{def:b2:regular-languages:regular}{regular}: $\{w_1, \dots, w_k\}$ is denoted by $w_1 \mid \cdots
\mid w_k$.

\begin{definition}[Extended regular expression]\label{def:b2:regular-languages:extended}
An \emph{extended regular expression}\index{extended regular expression} also uses these abbreviations: a class
\code{[abc]} for $a \mid b \mid c$ and a range \code{[a-z]} for every letter from \code{a} to \code{z}; $e^+$ for
$ee^*$; $e?$ for $e \mid \varepsilon$; and the counts $e\{m,n\}$ for $m$ copies of $e$ followed by $n - m$ copies of
$e?$, with $e\{m\} = e\{m,m\}$.
\end{definition}

Each abbreviation is an ordinary expression, so extended expressions denote \omterm{def:b2:regular-languages:regular}{regular languages} too; but they can be
much shorter: $(a \mid b)\{20\}$ stands for an expression with 40 letters. The parser of \code{regex.rs} reads them by
recursive descent, one function per priority level; the function for the postfix operators expands them on the spot:

\omcode{code/bachelor-2/22-regular-languages/listings/regex.rs}{86}{110}{Postfix operators: star, plus, option and
counts, each rewritten into the basic syntax}\label{lst:b2:regular-languages:post}

The syntax here is a small one. Practical pattern languages add a wildcard for any character, anchors for the start
and end of a line, and often \emph{back-references}, which ask a later part of the text to repeat what an earlier
group matched. Back-references go beyond \omterm{def:b2:regular-languages:regular}{regular languages} (\cref{sec:b2:regular-languages:limits}); the other
additions do not.

\section{Kleene's theorem}\label{sec:b2:regular-languages:kleene}

Expressions describe languages; automata recognise them. This section turns each into the other.

\begin{method}[Thompson's construction]\label{met:b2:regular-languages:thompson}
For each node of the expression's tree, build an NFA \emph{fragment} with \omterm{def:b2:finite-automata:epsilon}{spontaneous transitions}, one entry state $i$
and one exit state $o$, no transition into $i$ and none out of $o$, by induction on the tree
as drawn below:
\begin{itemize}
\item for $\emptyset$: two states and no transition; for $\varepsilon$: $i \xrightarrow{\varepsilon} o$; for a letter:
$i \xrightarrow{a} o$;
\item for $ef$: $i$ to the entry of $e$, the exit of $e$ to the entry of $f$, the exit of $f$ to $o$, all spontaneous;
\item for $e \mid f$: spontaneous moves from $i$ to both entries and from both exits to $o$;
\item for $e^*$: spontaneous moves from $i$ to the entry of $e$ and to $o$, and from the exit of $e$ back to its entry
and to $o$.
\end{itemize}
The fragment of the whole expression, with $i$ as start and $o$ as only \omterm{def:b2:finite-automata:dfa}{accepting state}, recognises $L(e)$. It has two
states per node, $2s$ for an expression of size $s$, and at most two transitions leave each state.
\end{method}

\begin{omfigure}
\begin{tikzpicture}[->, >={Stealth[length=1.6mm]}, every state/.style={minimum size=5mm, inner sep=0pt},
    font=\scriptsize, frag/.style={draw, rounded corners, minimum width=13mm, minimum height=6mm}, auto]
\begin{scope}
\node[state] (i) at (0, 0) {$i$}; \node[state] (o) at (1.5, 0) {$o$};
\path (i) edge node {$a$} (o);
\node at (0.75, -1.3) {letter $a$};
\end{scope}
\begin{scope}[xshift=2.6cm]
\node[state] (i) at (0, 0) {$i$}; \node[frag] (e) at (1.6, 0) {$e$}; \node[frag] (f) at (3.4, 0) {$f$};
\node[state] (o) at (5, 0) {$o$};
\path (i) edge node {$\varepsilon$} (e) (e) edge node {$\varepsilon$} (f) (f) edge node {$\varepsilon$} (o);
\node at (2.5, -1.3) {concatenation $ef$};
\end{scope}
\begin{scope}[yshift=-2.6cm]
\node[state] (i) at (0, 0) {$i$}; \node[frag] (e) at (1.8, 0.55) {$e$}; \node[frag] (f) at (1.8, -0.55) {$f$};
\node[state] (o) at (3.6, 0) {$o$};
\path (i) edge node {$\varepsilon$} (e) (i) edge node[below left] {$\varepsilon$} (f) (e) edge node {$\varepsilon$}
  (o) (f) edge node[below right] {$\varepsilon$} (o);
\node at (1.8, -1.4) {alternation $e \mid f$};
\end{scope}
\begin{scope}[xshift=5cm, yshift=-2.6cm]
\node[state] (i) at (0, 0) {$i$}; \node[frag] (e) at (1.8, 0) {$e$}; \node[state] (o) at (3.6, 0) {$o$};
\path (i) edge node {$\varepsilon$} (e) (e) edge node {$\varepsilon$} (o)
  (i) edge[bend right=40] node[below] {$\varepsilon$} (o)
  (e.north east) edge[bend right=60] node[above] {$\varepsilon$} (e.north west);
\node at (1.8, -1.4) {star $e^*$};
\end{scope}
\end{tikzpicture}
\par\smallskip{\small Thompson's fragments. A box is the fragment of a subexpression, entered on its left and left on
its right; every arrow but the letter's is a \omterm{def:b2:finite-automata:epsilon}{spontaneous transition}.}
\end{omfigure}

Correctness is a structural induction: a path from $i$ to $o$ reads exactly a word of the subexpression's language,
since the only way out of a fragment is through its exit, which no transition leaves inside it. Thompson's automaton
is easy to build but full of spontaneous moves. Glushkov's has none, and only one state per letter.

\begin{definition}[Linear expression, local language]\label{def:b2:regular-languages:local}
An expression is \emph{linear}\index{linear regular expression} if each letter occurs in it at most once. A language
$L$ is \emph{local}\index{local language} if there are sets $P, S \subseteq \Sigma$ and $F \subseteq \Sigma^2$ such
that $L \setminus \{\varepsilon\}$ is the set of nonempty words whose first letter is in $P$, whose last letter is in
$S$, and whose factors of length 2 are all in $F$.
\end{definition}

Whether a word belongs to a \omterm{def:b2:regular-languages:local}{local language} is decided by a window of two letters slid along it. Any expression is made
\omterm{def:b2:regular-languages:local}{linear} by numbering its letters from the left, its \emph{positions}: $(a \mid b)^*abb$ becomes $(a_1 \mid b_2)^* a_3
b_4 b_5$. Erasing the numbers maps the \omterm{def:b2:regular-languages:local}{linear} expression's language onto the original one.

\begin{definition}[Glushkov automaton]\label{def:b2:regular-languages:glushkov}
Let $e$ have positions $1, \dots, m$, the letter at position $j$ being $\ell(j)$, and let $e'$ be the \omterm{def:b2:regular-languages:local}{linear}
expression. Define $\mathrm{Null}$ (whether $\varepsilon \in L(e)$), $\mathrm{First}$ (the positions that can begin a
word of $L(e')$), $\mathrm{Last}$ (those that can end one), and $\mathrm{Follow}(j)$ (the positions that can come right
after $j$ in a word of $L(e')$). The \emph{Glushkov automaton}\index{Glushkov automaton} of $e$ has states $0, 1,
\dots, m$, start state 0, transitions $0 \xrightarrow{\ell(j)} j$ for $j \in \mathrm{First}$ and $i \xrightarrow{\ell(j)}
j$ for $j \in \mathrm{Follow}(i)$, and \omterm{def:b2:finite-automata:dfa}{accepting states} $\mathrm{Last}$, plus 0 if $\mathrm{Null}$.
\end{definition}

\begin{method}[The Berry--Sethi computation]\label{met:b2:regular-languages:berry-sethi}
Compute $\mathrm{Null}$, $\mathrm{First}$ and $\mathrm{Last}$ for every subexpression, bottom up: a position $j$ has
$(\mathrm{false}, \{j\}, \{j\})$; for $\varepsilon$, $(\mathrm{true}, \emptyset, \emptyset)$; an alternation takes the
\emph{or} and the unions; a concatenation $fg$ has $\mathrm{Null} = \mathrm{Null}(f) \wedge \mathrm{Null}(g)$,
$\mathrm{First} = \mathrm{First}(f)$, joined by $\mathrm{First}(g)$ if $f$ is nullable, and symmetrically for
$\mathrm{Last}$; a star is nullable with its body's sets. On the way, at each concatenation $fg$ add
$\mathrm{First}(g)$ to $\mathrm{Follow}(j)$ for every $j \in \mathrm{Last}(f)$, and at each star $f^*$ add
$\mathrm{First}(f)$ to $\mathrm{Follow}(j)$ for every $j \in \mathrm{Last}(f)$.
\end{method}

The walk of \code{regex.rs} is this method, one case per kind of node:

\omcode{code/bachelor-2/22-regular-languages/listings/regex.rs}{207}{241}{The Berry--Sethi computation: three values
per subexpression, and the follow sets filled at each concatenation and each
star}\label{lst:b2:regular-languages:walk}

\omoutput{code/bachelor-2/22-regular-languages/outputs/glushkov.txt}{The positions and sets of $(a \mid b)^*abb$, and
the automaton checked against the denotation.}\label{out:b2:regular-languages:glushkov}

\begin{omfigure}
\begin{tikzpicture}[->, >={Stealth[length=1.8mm]}, every state/.style={minimum size=7mm, inner sep=0pt},
    font=\scriptsize, auto]
\node[state, initial, initial text=] (0) at (0, 0) {0};
\node[state] (1) at (2.6, 1.3) {$a_1$};
\node[state] (2) at (2.6, -1.3) {$b_2$};
\node[state] (3) at (5.6, 0) {$a_3$};
\node[state] (4) at (7.6, 0) {$b_4$};
\node[state, accepting] (5) at (9.6, 0) {$b_5$};
\path (0) edge node {$a$} (1) (0) edge node[below left] {$b$} (2) (0) edge node[pos=0.3] {$a$} (3)
  (1) edge[loop above] node {$a$} () (2) edge[loop below] node {$b$} ()
  (1) edge[bend left=12] node {$b$} (2) (2) edge[bend left=12] node {$a$} (1)
  (1) edge node {$a$} (3) (2) edge node[below right] {$a$} (3)
  (3) edge node {$b$} (4) (4) edge node {$b$} (5);
\end{tikzpicture}
\par\smallskip{\small The \omterm{def:b2:regular-languages:glushkov}{Glushkov automaton} of $(a \mid b)^*abb$: one state per position, and every transition into
a position carries that position's letter. It is nondeterministic: from 0, $a_1$ and $a_3$ both read
$a$.}
\end{omfigure}

\begin{proposition}[Glushkov's automaton]\label{prop:b2:regular-languages:glushkov}
The \omterm{def:b2:regular-languages:glushkov}{Glushkov automaton} of an expression with $m$ positions has $m + 1$ states, no \omterm{def:b2:finite-automata:epsilon}{spontaneous transition}, and at most
$m(m + 1)$ transitions, and it recognises $L(e)$.
\end{proposition}

\begin{proof}
The counts come from the definition. For the language, first show by structural induction that $L(e')$ is \omterm{def:b2:regular-languages:local}{local} with
$P = \mathrm{First}$, $S = \mathrm{Last}$ and $F = \{ij : j \in \mathrm{Follow}(i)\}$, and that the sets computed by
\cref{met:b2:regular-languages:berry-sethi} are these. For a position, both are clear. For $f \mid g$, the positions of
$f$ and $g$ are disjoint, no follow pair joins them, so a nonempty word passing the window test uses the positions of
one side only and belongs to it. For $fg$, a word passing the test splits at the only pair $ij$ with $i$ in $f$ and $j$
in $g$, which the method added exactly for $i \in \mathrm{Last}(f)$, $j \in \mathrm{First}(g)$; the two parts pass the
tests of $f$ and $g$, or one part is empty and its side nullable. For $f^*$, a word passing the test splits at the
pairs $ij$ with $i \in \mathrm{Last}(f)$, $j \in \mathrm{First}(f)$ added by the star, into words of $L(f')$; and
conversely every word of these languages passes the test. Finally, a run of the automaton on $w = a_1 \cdots a_k$
from 0 visits positions $j_1, \dots, j_k$ with $\ell(j_t) = a_t$, $j_1 \in \mathrm{First}$ and each $j_{t+1} \in
\mathrm{Follow}(j_t)$; it accepts when $j_k \in \mathrm{Last}$ (or $k = 0$ and $\mathrm{Null}$). So the automaton
accepts $w$ exactly when $w$ is the image of a word of the \omterm{def:b2:regular-languages:local}{local language} $L(e')$, that is, $w \in L(e)$.
\end{proof}

Thompson's automaton for $(a \mid b)^*abb$ has 20 states; Glushkov's has 6:

\omoutput{code/bachelor-2/22-regular-languages/outputs/thompson.txt}{Thompson's construction, two states per node of
the tree, against Glushkov's $m + 1$, and the two automata compared on every word of length at most
7.}\label{out:b2:regular-languages:thompson}

The way back, from an automaton to an expression, labels transitions with expressions and removes the states one by
one.

\begin{method}[State elimination]\label{met:b2:regular-languages:elimination}
Given an automaton with states $1, \dots, n$, add a new start $S$ with $S \xrightarrow{\varepsilon}$ the old start and
a new end $E$ with $q \xrightarrow{\varepsilon} E$ for every accepting $q$. Label each pair $(i, j)$ by an expression
$r_{ij}$: the alternation of the letters from $i$ to $j$, or $\emptyset$. Then, for $k = 1, \dots, n$, remove state
$k$: for every pair $(i, j)$ of remaining states, replace
\[
r_{ij} \quad\text{by}\quad r_{ij} \mid r_{ik}\, r_{kk}^*\, r_{kj}.
\]
At the end only $S$ and $E$ remain, and $r_{SE}$ denotes the language of the automaton.
\end{method}

\begin{omfigure}
\begin{tikzpicture}[->, >={Stealth[length=1.8mm]}, every state/.style={minimum size=6mm, inner sep=0pt},
    font=\scriptsize, auto]
\begin{scope}
\node[state] (S) at (0, 0) {$S$}; \node[state] (0) at (1.8, 0) {0}; \node[state] (1) at (3.8, 0) {1};
\node[state] (E) at (5.6, 0) {$E$};
\path (S) edge node {$\varepsilon$} (0) (0) edge[loop above] node {$a$} () (0) edge[bend left=20] node {$b$} (1)
  (1) edge[bend left=20] node {$a$} (0) (1) edge[loop above] node {$b$} () (1) edge node {$\varepsilon$} (E);
\node[anchor=west] at (6.2, 0) {the automaton of ``ends with $b$''};
\end{scope}
\begin{scope}[yshift=-2cm]
\node[state] (S) at (0, 0) {$S$}; \node[state] (1) at (3.8, 0) {1}; \node[state] (E) at (5.6, 0) {$E$};
\path (S) edge node {$a^*b$} (1) (1) edge[loop above] node {$b \mid aa^*b$} () (1) edge node {$\varepsilon$} (E);
\node[anchor=west] at (6.2, 0) {state 0 removed};
\end{scope}
\begin{scope}[yshift=-3.4cm]
\node[state] (S) at (0, 0) {$S$}; \node[state] (E) at (5.6, 0) {$E$};
\path (S) edge node {$a^*b(b \mid aa^*b)^*$} (E);
\node[anchor=west] at (6.2, 0) {state 1 removed};
\end{scope}
\end{tikzpicture}
\par\smallskip{\small State elimination in three steps. Removing 0 reroutes $S \to 0 \to 1$ into $a^*b$ (the loop
$a$ on 0 becomes $a^*$) and $1 \to 0 \to 1$ into $aa^*b$, joined to the loop $b$.}
\end{omfigure}

The loop in \code{dfa.rs}, the second library of the chapter, is the formula of the method; the smart \omterm{def:b2:cpp-classes-raii:constructor}{constructors}
\code{cat}, \code{alt} and \code{star} of \code{regex.rs} drop the $\emptyset$ and $\varepsilon$ that would otherwise
fill the result:

\omcode{code/bachelor-2/22-regular-languages/listings/dfa.rs}{115}{140}{State elimination on a complete
DFA}\label{lst:b2:regular-languages:eliminate}

\omoutput{code/bachelor-2/22-regular-languages/outputs/eliminate.txt}{The three-state automaton of the binary multiples
of 3 (\cref{ch:b2:finite-automata}), turned into an expression and checked.}\label{out:b2:regular-languages:eliminate}

\begin{theorem}[Kleene's theorem]\label{thm:b2:regular-languages:kleene}
A language is \omterm{def:b2:regular-languages:regular}{regular} if and only if it is \omterm{def:b2:finite-automata:recognisable}{recognisable} by a finite automaton.
\end{theorem}

\begin{proof}
If $L = L(e)$, the \omterm{def:b2:regular-languages:glushkov}{Glushkov automaton} recognises it (\cref{prop:b2:regular-languages:glushkov}). Conversely, take an
automaton and run \cref{met:b2:regular-languages:elimination}. By induction on $k$: after states $1, \dots, k$ are
removed, $r_{ij}$ denotes the set of words that lead from $i$ to $j$ through intermediate states among $1, \dots, k$
only. True for $k = 0$, where there is no intermediate state. A path from $i$ to $j$ through $1, \dots, k$ either avoids
$k$, and is counted by the old $r_{ij}$, or cuts at its visits to $k$ into a piece from $i$ to $k$, any number of loops
from $k$ to $k$, and a piece from $k$ to $j$, each avoiding $k$ inside: that is $r_{ik} r_{kk}^* r_{kj}$. At the end,
$r_{SE}$ denotes the words leading from $S$ to $E$, which are the accepted words.
\end{proof}

With \cref{ch:b2:finite-automata}, \omterm{def:b2:regular-languages:regular}{regular languages} are therefore closed under complement and intersection too,
though expressions have no operator for either: the way goes through automata.

\begin{history}
% ledger: b2:kleene-1956
% ledger: b2:glushkov-1961
% ledger: b2:thompson-1968
% ledger: b2:berry-sethi-1986
Stephen Kleene introduced \omterm{def:b2:regular-languages:regex}{regular expressions} and proved the theorem in his paper of 1956 on nerve nets and finite
automata. Victor Glushkov's position automaton appeared in his 1961 survey of the theory of automata. Ken Thompson
published his construction in 1968, in a short paper on searching text for a \omterm{def:b2:regular-languages:regex}{regular expression}, where the NFA is
simulated on all its states at once. Gérard Berry and Ravi Sethi showed in
1986 how to compute the positions' sets on the expression directly, as in \cref{met:b2:regular-languages:berry-sethi}.
\end{history}

\section{The minimal automaton}\label{sec:b2:regular-languages:minimal}

Many automata recognise the same language. What a DFA must remember after reading $u$ is what the rest of the word
needs to know.

\begin{definition}[Residual, Myhill--Nerode equivalence]\label{def:b2:regular-languages:residual}
The \emph{residual}\index{residual (language)} (or left quotient) of a language $L$ by a word $u$ is $u^{-1}L = \{v :
uv \in L\}$. Two words are \emph{equivalent for $L$}\index{Myhill--Nerode equivalence}, $u \sim_L v$, if $u^{-1}L =
v^{-1}L$: no continuation $x$ tells them apart, $ux \in L \iff vx \in L$.
\end{definition}

For $L = (ab \mid b)^*$: $\varepsilon^{-1}L = L$, $a^{-1}L = bL$, $(aa)^{-1}L = \emptyset$, and $b^{-1}L = (ab)^{-1}L =
L$. Every word has one of these three \omterm{def:b2:regular-languages:residual}{residuals}.

\begin{theorem}[Myhill--Nerode]\label{thm:b2:regular-languages:myhill-nerode}
A language $L$ is \omterm{def:b2:finite-automata:recognisable}{recognisable} if and only if it has finitely many \omterm{def:b2:regular-languages:residual}{residuals}. Every complete DFA for $L$ has at least as
many \omterm{def:b2:finite-automata:trim}{accessible states} as $L$ has \omterm{def:b2:regular-languages:residual}{residuals}, and the \emph{\omterm{def:b2:regular-languages:residual}{residual} automaton} reaches the bound: its states are the
\omterm{def:b2:regular-languages:residual}{residuals}, its start is $L$, its transitions are $R \xrightarrow{a} a^{-1}R$, and its \omterm{def:b2:finite-automata:dfa}{accepting states} are the
\omterm{def:b2:regular-languages:residual}{residuals} that contain $\varepsilon$.
\end{theorem}

\begin{proof}
Let $A$ be a complete DFA for $L$ and $L_q$ the language it accepts from state $q$. Then $u^{-1}L = L_{\delta^*(q_0,
u)}$: $uv$ is accepted exactly when $v$ leads $\delta^*(q_0, u)$ to an \omterm{def:b2:finite-automata:dfa}{accepting state}
(\cref{prop:b2:finite-automata:extended}). So each \omterm{def:b2:regular-languages:residual}{residual} is $L_q$ for some accessible $q$, and there are at most as
many \omterm{def:b2:regular-languages:residual}{residuals} as \omterm{def:b2:finite-automata:trim}{accessible states}. Conversely, suppose there are finitely many \omterm{def:b2:regular-languages:residual}{residuals}. Since $a^{-1}(u^{-1}L) =
(ua)^{-1}L$, the transitions of the \omterm{def:b2:regular-languages:residual}{residual} automaton lead from a \omterm{def:b2:regular-languages:residual}{residual} to a \omterm{def:b2:regular-languages:residual}{residual}, and by induction on $|w|$ its
run on $w$ from $L$ ends in $w^{-1}L$. It accepts $w$ when $\varepsilon \in w^{-1}L$, that is, when $w \in L$.
\end{proof}

\begin{definition}[Minimal automaton]\label{def:b2:regular-languages:minimal}
A complete DFA for $L$ is \emph{minimal}\index{minimal automaton} if no complete DFA for $L$ has fewer states. By the
theorem, the \omterm{def:b2:regular-languages:residual}{residual} automaton is minimal; it is called \emph{the} minimal automaton of $L$.
\end{definition}

\begin{corollary}[Uniqueness]\label{cor:b2:regular-languages:unique-minimal}
Every \omterm{def:b2:regular-languages:minimal}{minimal} complete DFA for $L$ equals the \omterm{def:b2:regular-languages:residual}{residual} automaton up to the names of its states.
\end{corollary}

\begin{proof}
In a \omterm{def:b2:regular-languages:minimal}{minimal} DFA $A$ every state is accessible, or removing it would leave a smaller complete DFA. The map $q \mapsto
L_q$ sends each state to a \omterm{def:b2:regular-languages:residual}{residual}, reaches them all, and so is a bijection, since both sets have the same size. It
respects the start ($L_{q_0} = L$), the \omterm{def:b2:finite-automata:dfa}{accepting states} ($q \in F$ iff $\varepsilon \in L_q$) and the transitions
($L_{\delta(q, a)} = a^{-1}L_q$).
\end{proof}

Two states $p$ and $q$ of a DFA can be merged when $L_p = L_q$. Moore's method finds these pairs by looking further and
further ahead.

\begin{method}[Moore's minimisation]\label{met:b2:regular-languages:moore}
Start from an accessible complete DFA with $n$ states. Round 0 splits the states into two classes, accepting or not.
Round $k + 1$ puts $p$ and $q$ in the same class when they were in the same class at round $k$ and, for every letter
$a$, so were $\delta(p, a)$ and $\delta(q, a)$. Stop at the first round that splits no class. Merging each final class
into one state gives the \omterm{def:b2:regular-languages:minimal}{minimal automaton}.
\end{method}

\begin{proof}[Correctness]
By induction on $k$, $p$ and $q$ share a class at round $k$ exactly when no word of length at most $k$ is accepted
from one and not from the other. A round that splits nothing gives the same partition at every later round, since a
round depends only on the previous one; so the final classes are those of $L_p = L_q$. Merging them gives a well
defined DFA, since equivalent states have equivalent successors, with one state per \omterm{def:b2:regular-languages:residual}{residual} of $L$: the \omterm{def:b2:regular-languages:minimal}{minimal
automaton}. Each round but the last adds a class, and there are at most $n$, so there are at most $n$ rounds, each in
time proportional to $n|\Sigma|$.
\end{proof}

\omcode{code/bachelor-2/22-regular-languages/listings/dfa.rs}{53}{80}{Moore's rounds: a state's new class is named by
its old class and its successors' classes}\label{lst:b2:regular-languages:moore}

The determinised \omterm{def:b2:regular-languages:glushkov}{Glushkov automaton} of ``starts with $a$, or starts and ends with $b$'' has 7 states; two rounds
leave 4 classes:

\omoutput{code/bachelor-2/22-regular-languages/outputs/minimise.txt}{The DFA of $a(a \mid b)^* \mid b(a \mid b)^*b
\mid b$, Moore's classes round by round, and the result.}\label{out:b2:regular-languages:minimise}

\begin{omfigure}
\begin{tikzpicture}[->, >={Stealth[length=1.8mm]}, every state/.style={minimum size=9mm, inner sep=0pt},
    font=\scriptsize, auto]
\node[state, initial, initial text=] (A) at (0, 0) {$\{0\}$};
\node[state, accepting] (B) at (2.8, 1) {$\{1, 3, 4\}$};
\node[state, accepting] (C) at (2.8, -1) {$\{2, 6\}$};
\node[state] (D) at (5.8, -1) {$\{5\}$};
\path (A) edge node {$a$} (B) (A) edge node[below left] {$b$} (C) (B) edge[loop right] node {$a, b$} ()
  (C) edge[loop below] node {$b$} () (C) edge[bend left=15] node {$a$} (D) (D) edge[bend left=15] node {$b$} (C)
  (D) edge[loop right] node {$a$} ();
\end{tikzpicture}
\par\smallskip{\small The \omterm{def:b2:regular-languages:minimal}{minimal automaton}: each state is a class of \cref{out:b2:regular-languages:minimise}. In
$\{2, 6\}$ the last letter read is a $b$ after a first $b$; in $\{5\}$ it is an $a$.}
\end{omfigure}

Moore's method can need $n$ rounds, on an automaton shaped as a chain, and then costs time proportional to
$n^2|\Sigma|$. John Hopcroft's algorithm reaches $n |\Sigma| \log n$: it keeps a list of \emph{splitters}, a class and a
letter, splits every class by ``enters the splitter on that letter or not'', and when a class splits in two puts only
the smaller half on the list. A state then joins a splitter at most $\log_2 n$ times, since each time its class at
least halves (\cref{exo:b2:regular-languages:11}).

\section{The limits of regular languages}\label{sec:b2:regular-languages:limits}

A finite automaton has a finite memory. A language that needs to count without bound is not \omterm{def:b2:regular-languages:regular}{regular}.

\begin{theorem}[Pumping lemma]\label{thm:b2:regular-languages:pumping}
If $L$ is recognised by a DFA with $p$ states, every word $w \in L$ with $|w| \ge p$ can be written $w = xyz$ with
$|xy| \le p$ and $y \ne \varepsilon$, so that $xy^kz \in L$ for every $k \ge 0$.
\end{theorem}

\begin{proof}
The run on the first $p$ letters of $w$ visits $p + 1$ states, so one state $q$ appears twice, after reading $x$ and
after reading $xy$, with $|xy| \le p$ and $y$ nonempty. Reading $y$ leads $q$ back to $q$, so it can be read any
number of times, zero included, and the run on $xy^kz$ ends where the run on $w$ ends, in an \omterm{def:b2:finite-automata:dfa}{accepting state}.
\end{proof}

\begin{proposition}[Three languages that are not regular]\label{prop:b2:regular-languages:nonregular}
The languages $\{a^n b^n : n \ge 0\}$, the balanced bracket words over $\{(, )\}$, and the palindromes over $\{a, b\}$
are not \omterm{def:b2:regular-languages:regular}{regular}.
\end{proposition}

\begin{proof}
For $\{a^nb^n\}$, by the pumping lemma: if a DFA with $p$ states recognised it, $w = a^pb^p$ would split with $y =
a^j$, $j \ge 1$, since $|xy| \le p$; but $xy^2z = a^{p+j}b^p$ is not in the language. Also by
\cref{thm:b2:regular-languages:myhill-nerode}: for $i \ne j$, $b^i$ is in the \omterm{def:b2:regular-languages:residual}{residual} by $a^i$ and not in the \omterm{def:b2:regular-languages:residual}{residual}
by $a^j$, so the \omterm{def:b2:regular-languages:residual}{residuals} are infinitely many. Balanced brackets: their intersection with the \omterm{def:b2:regular-languages:regular}{regular language}
$(^*\,)^*$ is $\{(^n )^n\}$, which is $\{a^nb^n\}$ with other letters; if they were \omterm{def:b2:regular-languages:regular}{regular}, so would be the
intersection. Palindromes: for $i \ne j$, $ba^i$ is in the \omterm{def:b2:regular-languages:residual}{residual} by $a^ib$ and not in the \omterm{def:b2:regular-languages:residual}{residual} by $a^jb$.
\end{proof}

The pumping lemma gives only a necessary condition: some languages that are not \omterm{def:b2:regular-languages:regular}{regular} can be pumped. The \omterm{def:b2:regular-languages:residual}{residuals}
give the exact answer, and with them a proof is a matter of finding infinitely many words that pairwise differ in what
must follow. A back-reference that repeats a group asks for the language $\{ww\}$, which has infinitely many \omterm{def:b2:regular-languages:residual}{residuals}
(by $a^ib$, as for palindromes); this is why pattern engines with back-references cannot be finite automata. Balanced
brackets and arithmetic expressions need the grammars of \cref{ch:b2:context-free-languages}.

% ledger: b2:nerode-1958
\begin{history}
The equivalence and the theorem carry the names of John Myhill and Anil Nerode; Nerode's paper appeared in 1958. The
pumping lemma for \omterm{def:b2:regular-languages:regular}{regular languages} is a direct form of the pigeonhole argument used in its proof.
\end{history}

\section{Deciding questions about automata}\label{sec:b2:regular-languages:decisions}

Because an automaton is a finite object, many questions about the infinite language it recognises can be answered by a
program that always stops.

\begin{proposition}[Questions and their costs]\label{prop:b2:regular-languages:decisions}
For complete DFAs $A$ and $B$ with $n_A$ and $n_B$ states over $\Sigma$:
\begin{itemize}
\item \emph{membership} of a word $w$: time proportional to $|w|$;
\item \emph{emptiness}, $L(A) = \emptyset$: no \omterm{def:b2:finite-automata:dfa}{accepting state} is accessible, a graph search in time proportional to
$n_A|\Sigma|$;
\item \emph{finiteness}: $L(A)$ is infinite exactly when the \omterm{def:b2:finite-automata:trim}{trim automaton} has a cycle, also a graph search;
\item \emph{inclusion}, $L(A) \subseteq L(B)$: $L(A) \cap (\Sigma^* \setminus L(B))$ is empty, a search of the product
in time proportional to $n_A n_B |\Sigma|$;
\item \emph{equivalence}: no accessible pair $(p, q)$ of the product has $p$ accepting and $q$ not, or the reverse; a
breadth-first search finds a shortest word on which they differ, if there is one.
\end{itemize}
\end{proposition}

\begin{proof}
Membership is the run. Emptiness: a word is accepted exactly when its run reaches an \omterm{def:b2:finite-automata:dfa}{accepting state}, so $L(A)$ is
nonempty exactly when an \omterm{def:b2:finite-automata:dfa}{accepting state} is accessible. Finiteness: in the \omterm{def:b2:finite-automata:trim}{trim automaton} every path extends into an
accepted word; a cycle gives infinitely many paths, and without a cycle every path has fewer than $n_A$ transitions, so
the accepted words are shorter than $n_A$ and finitely many. Inclusion is emptiness of the product for intersection
with the complement (\cref{prop:b2:finite-automata:closure}), and equivalence is two inclusions, read on the same
product: the pair after $w$ is $(\delta_A^*(q_A, w), \delta_B^*(q_B, w))$, and $w$ separates the languages exactly when
one of the two states accepts and the other does not.
\end{proof}

From expressions, these questions go through automata. The \omterm{def:b2:regular-languages:glushkov}{Glushkov automaton} is built in time proportional to $m^2$
and decides membership by simulation in at most $(m + 1)|w|$ steps; but inclusion and equivalence need complete DFAs,
and determinisation may cost exponentially many states (\cref{thm:b2:finite-automata:blowup}). For short patterns it is
fast:

\omoutput{code/bachelor-2/22-regular-languages/outputs/equivalent.txt}{Five pairs of expressions, compared through
their determinised Glushkov automata, with a shortest witness when they
differ.}\label{out:b2:regular-languages:equivalent}

The witness \code{"bb"} is in the language of $(a \mid b)^*bb$ but does not end with $abb$. The checker itself is
\cref{exo:b2:regular-languages:9}. By \cref{cor:b2:regular-languages:unique-minimal}, a second way to compare is to
minimise both automata and compare them state by state.

\section{A pattern engine that cannot freeze}\label{sec:b2:regular-languages:engine}

The engine of \code{regex.rs} parses the pattern into a tree, computes Glushkov's sets, and simulates the automaton
on sets of states, as \cref{ch:b2:finite-automata} simulated NFAs. With at most 127 positions, a set of states is the
bits of one 128-bit integer; for each letter, the next set is the union of the successors of the current states, kept
only at positions holding that letter:

\omcode{code/bachelor-2/22-regular-languages/listings/regex.rs}{251}{278}{The matcher: one OR per current state,
then a mask for the letter}\label{lst:b2:regular-languages:matcher}

\begin{proposition}[Linear matching]\label{prop:b2:regular-languages:linear}
On an expression with $m$ positions and a text of $n$ letters, the matcher performs at most $(m + 1)n$ unions, whatever
the expression and the text.
\end{proposition}

\begin{proof}
Each letter costs one union per state of the current set, and the set holds at most the $m + 1$ states. The set after
$k$ letters is the set of states the \omterm{def:b2:regular-languages:glushkov}{Glushkov automaton} can reach on the first $k$ letters, by induction on $k$, so the
answer is that of the automaton, which is correct by \cref{prop:b2:regular-languages:glushkov}.
\end{proof}

On $(a \mid a)^*b$, with positions $a_1$, $a_2$, $b_3$, the set after each $a$ is $\{1, 2\}$, and every letter costs
two unions:

\omoutput{code/bachelor-2/22-regular-languages/outputs/engine_compare.txt}{The matcher on $(a \mid a)^*b$ and $a^n$,
on the laptop.}\label{out:b2:regular-languages:engine}

\begin{omfigure}
\begin{tikzpicture}
\begin{axis}[omrecord, width=0.72\linewidth, height=5cm, xlabel={number $n$ of letters $a$}, ylabel={time (s)},
    ymode=log, xmin=15.5, xmax=27.5, ymin=1e-8, ymax=20, legend style={font=\scriptsize, at={(0.97,0.5)}, anchor=east}]
  \addplot[omDef, thick, mark=*, mark size=1pt] table[x=n, y=seconds]
    {code/bachelor-2/22-regular-languages/figures/backtrack.dat};
  \addlegendentry{Python's backtracking \texttt{re}}
  \addplot[omThm, thick, dashed, domain=16:27] {7.34e-9 * x};
  \addlegendentry{Glushkov matcher, 7.3 ns a letter}
\end{axis}
\end{tikzpicture}
\par\smallskip{\small Time to answer ``no match'' for $(a \mid a)^*b$ on $a^n$, on a logarithmic scale: the
backtracking engine's points (\cref{out:b2:regular-languages:hook}) climb on a straight line, doubling with each $a$;
the matcher's line is its measured cost per letter on ten million letters (\cref{out:b2:regular-languages:engine}),
times $n$.}
\end{omfigure}

A backtracking engine tries one way to match at a time: on $a^n$, each $a$ can be matched by the left or the right
branch of $(a \mid a)$, and when the missing $b$ makes a choice fail, the engine goes back to the last choice and tries
the other. The $2^n$ ways to read $a^n$ all fail, and it tries them all. The automaton keeps every way at once, in a
set that never holds more than $m + 1$ states. Backtracking has reasons: it reports which part of the text each group
matched, in the order a programmer expects, and it supports back-references, which no automaton can. Engines built on
automata, in the line of Thompson's, give up back-references and in exchange answer in \omterm{def:b2:regular-languages:local}{linear} time.

\begin{dated}{2026}
% ledger: b2:python-re-backtracking
% ledger: b2:posix-ere
Python's \code{re} module is a backtracking engine. Since version 3.11 it offers possessive quantifiers such as
\code{a*+} and \omterm{def:b2:threads-and-synchronisation:atomic}{atomic} groups, which throw away the backtracking points inside them and can stop a blow-up like the
hook's when the pattern's author sees it coming. The POSIX standard defines basic and \omterm{def:b2:regular-languages:extended}{extended regular expressions} for
tools such as \code{grep}, and asks, when several matches start at the same point, for the longest.
\end{dated}

On the town's 24 street names, the engine answers patterns that expand to a hundred positions, since a range such as
\code{[A-Z]} is an alternation of 26 letters:

\omcode{code/bachelor-2/22-regular-languages/listings/streets.rs}{1}{21}{Five patterns over the street
names}\label{lst:b2:regular-languages:streets}

\omoutput{code/bachelor-2/22-regular-languages/outputs/streets.txt}{The number of positions of each pattern, the
number of names it matches, and the first two.}\label{out:b2:regular-languages:streets}

\section{Exercises}\label{sec:b2:regular-languages:exercises}

\begin{exercise}[$\star$]\label{exo:b2:regular-languages:1}
Which of the words $\varepsilon$, $a$, $aa$, $bab$, $abba$, $babab$, $aab$ are in $L((ab \mid b)^*a?)$? Describe the
language in words, then check with the engine.
\end{exercise}

\begin{exercise}[$\star$]\label{exo:b2:regular-languages:2}
A tool checks variable names with the pattern \code{[a-z0-9_]+}. Find the names it wrongly accepts, and write a
pattern that requires a first character that is not a digit.
\end{exercise}

\begin{exercise}[$\star$]\label{exo:b2:regular-languages:3}
Write an extended expression for the decimal numbers of \cref{ch:b2:finite-automata}'s hook (an optional sign,
digits, an optional point followed by digits, an optional exponent). Check it against the hook's automaton on every
string of length at most 6 over the characters \code{0 1 + - . e E}; in this chapter's syntax, \code{.} is a plain
letter.
\end{exercise}

\begin{exercise}[$\star$]\label{exo:b2:regular-languages:4}
Let $c_n$ be the number of words of length $n$ in $L((ab \mid b)^*)$. Find a recurrence for $c_n$, and check it on
$c_0, \dots, c_{12}$.
\end{exercise}

\begin{exercise}[$\star\star$]\label{exo:b2:regular-languages:5}
Number the positions of $a(ba \mid b)^*a$ and compute $\mathrm{Null}$, $\mathrm{First}$, $\mathrm{Last}$ and every
$\mathrm{Follow}$ set by \cref{met:b2:regular-languages:berry-sethi}. Draw the \omterm{def:b2:regular-languages:glushkov}{Glushkov automaton}.
\end{exercise}

\begin{exercise}[$\star\star$]\label{exo:b2:regular-languages:6}
A program minimises the incomplete DFA of $\{ab, ba\}$ (states 0, 1 after $a$, 2 after $b$, 3 accepting) by refining
classes on the classes of the successors that exist, without completing it first. Which states does it merge, and
which words does the merged automaton wrongly accept? How must the program be fixed?
\end{exercise}

\begin{exercise}[$\star\star$]\label{exo:b2:regular-languages:7}
Apply \cref{met:b2:regular-languages:elimination} by hand to the two-state DFA of ``an even number of $a$'' over
$\{a, b\}$, removing the start state first. Check the expression on every word of length at most 10.
\end{exercise}

\begin{exercise}[$\star\star$]\label{exo:b2:regular-languages:8}
Prove that $\{a^n b^m : n \le m\}$ is not \omterm{def:b2:regular-languages:regular}{regular}, once with the pumping lemma and once with \omterm{def:b2:regular-languages:residual}{residuals}.
\end{exercise}

\begin{exercise}[$\star\star$]\label{exo:b2:regular-languages:9}
Write a function that takes two complete DFAs over the same \omterm{def:b2:finite-automata:word}{alphabet} and returns either ``equivalent'' or a shortest
word on which they differ. Use it to compare expressions through their Glushkov automata.
\end{exercise}

\begin{exercise}[$\star\star\star$]\label{exo:b2:regular-languages:10}
Show that $\sim_L$ is a \emph{right congruence}: $u \sim_L v$ implies $ua \sim_L va$ for every letter $a$. Deduce that
if $L$ is \omterm{def:b2:regular-languages:regular}{regular}, the classes of $\sim_L$ are in bijection with the states of the \omterm{def:b2:regular-languages:minimal}{minimal automaton}, $u \sim_L v$
holding exactly when $u$ and $v$ lead it to the same state.
\end{exercise}

\begin{exercise}[$\star\star\star$]\label{exo:b2:regular-languages:11}
Implement Hopcroft's algorithm: keep a list of splitters (a class and a letter); for a splitter $(B, a)$, split each
class into the states that enter $B$ on $a$ and the others; when a class splits, put both halves on the list if it
was there, and otherwise only the smaller half. Check that it finds Moore's classes, and compare the two on large
automata. Why is the gain so small here?
\end{exercise}

\begin{exercise}[$\star\star\star$]\label{exo:b2:regular-languages:12}
Prove that $\{w \in \{a, b\}^* : w \text{ has as many } a \text{ as } b\}$ is not \omterm{def:b2:regular-languages:regular}{regular}, once with \omterm{def:b2:regular-languages:residual}{residuals} and
once by intersecting it with $a^*b^*$.
\end{exercise}

\section{Problem: A search box that never freezes}\label{sec:b2:regular-languages:problem}

\begin{problem}[{Weekend problem --- A search box that never freezes}]\label{pb:b2:regular-languages:1}
The town's website lets visitors search the street names with patterns. The first version used a backtracking engine,
and a visitor froze it. This problem builds the replacement. The 24 names are Bridge Road, Bridge Way, Castle Road,
Castle Row, Castle Way, Church Lane, Church Way, Garden Lane, Garden Road, Market Road, Market Row, Mill Avenue, Mill
Row, Mill Street, Orchard Lane, Orchard Row, Orchard Way, River Avenue, River Lane, School Avenue, School Lane, School
Way, Station Row and Willow Lane.

\medskip\noindent\textbf{Part I --- Patterns over the street names.}\par\nopagebreak
\begin{enumerate}
\item How many names end in Road, in Row and in Way? Compare with \cref{out:b2:regular-languages:streets}.
\item Why does \code{[A-Z][a-z]* Road} have 57 positions?
\item Write a pattern for the names whose second word is Lane or Avenue. How many names does it match?
\item List the 8 names matched by \code{[A-Z][a-z]*(e|n) [A-Z][a-z]*}.
\item Show that any finite set of names is a \omterm{def:b2:regular-languages:regular}{regular language}.
\item A search box looks for the pattern anywhere in a name. Which expression, built from $e$, denotes the names that
contain a match of $e$?
\end{enumerate}

\medskip\noindent\textbf{Part II --- The Glushkov engine.}\par\nopagebreak
Let $e = (ab \mid b)^*c$ over $\{a, b, c\}$.
\begin{enumerate}[resume, beginpenalty=10000]
\item Number the positions of $e$ and give $\mathrm{Null}$, $\mathrm{First}$, $\mathrm{Last}$ and the follow sets.
\item How many states and transitions does its \omterm{def:b2:regular-languages:glushkov}{Glushkov automaton} have?
\item Give the matcher's set of states after each letter of $babbc$. Is the word accepted, and how many unions were
made?
\item Show that after reading a letter, every state in the set is a position holding that letter.
\item At most how many unions does the matcher make on a name of 12 letters with the pattern \code{[A-Z][a-z]* Road}?
\item Why does the matcher refuse patterns with more than 127 positions, and how would you lift the limit?
\end{enumerate}

\medskip\noindent\textbf{Part III --- One automaton per language.}\par\nopagebreak
\begin{enumerate}[resume, beginpenalty=10000]
\item Determinise the \omterm{def:b2:regular-languages:glushkov}{Glushkov automaton} of $e$ over $\{a, b, c\}$, with the empty set as sink. How many states?
\item Apply Moore's method to it. How many rounds, and how many classes at the end?
\item List the \omterm{def:b2:regular-languages:residual}{residuals} of $L(e)$, and match them with the final classes.
\item Show that $(ab \mid b)^*c$ and $b^*(abb^*)^*c$ are equivalent.
\item Find a shortest word on which $(ab \mid b)^*c$ and $(a \mid b)^*bc$ differ.
\item The site stores compiled patterns in a cache. Why can two patterns that denote the same language share one
entry, and how can a key be computed for it?
\end{enumerate}

\medskip\noindent\textbf{Part IV --- Backtracking against the automaton.}\par\nopagebreak
\begin{enumerate}[resume, beginpenalty=10000]
\item In how many ways can $a^n$ be read as a sequence of matches of $(a \mid a)$?
\item Why does a backtracking engine try all of them on $a^n$? Relate this to the ratios of
\cref{out:b2:regular-languages:hook}.
\item From the hook's last line, estimate how long the old site would take for 40 letters.
\item How many unions does the matcher make on $a^n$ with $(a \mid a)^*b$? Check with
\cref{out:b2:regular-languages:engine}.
\item How long does the matcher take for 40 letters, and for ten million?
\item State the result: from exponential to \omterm{def:b2:regular-languages:local}{linear}, backtracking's time doubling with each $a$, measured, against at
most $(m + 1)n$ steps proved for the automaton.
\end{enumerate}
\end{problem}
